\section{位置と速度}\label{節:位置と速度}
一直線上を動く物体を考えよう.直線上に原点を定め,右向きを正とする.時刻$t$における位置を$s(t)$と書けば,運動は1つの関数で表される.位置が$3$から$8$へ変わったなら,位置の変化は$8-3=5$である.逆に$8$から$3$へ変わったなら$3-8=-5$となる.負号によって,どちらへ動いたかも区別できる.

速度が一定なら,位置の変化は「速度$\times$時間」で求められる.速度が毎秒$2$メートルなら,\,$3$秒間の位置の変化は$6$メートルである.では,途中で速度が変わる場合はどうだろうか.どの時刻の速度を掛けるかで答えが変わってしまうし,そもそも速度が変わり続けるときの「時刻$2$の速度」とは何を指すのかもはっきりしない.まずは位置から速度を求めることを考えよう.

時刻$t$から$t+h$までの位置の変化は$s(t+h)-s(t)$である.これを経過した時間$h$で割った
\[
 \frac{s(t+h)-s(t)}h
\]
を,この区間の\textbf{平均の速度}という.位置の単位がメートル,時間の単位が秒なら,速度の単位はメートル毎秒となる.

ここで1つ断りをいれておこう.平均の速度は$h$で割って作るので,\,$h\ne0$でなければ意味をもたない.以後,このような分数では断りなしに$h\ne0$とする.また$h<0$のときは,時刻$t+h$から$t$までの平均の速度を表している.

例えば$s(t)=t^2$として,時刻$2$の前後で実験してみよう.\,$h$をいろいろに変えて平均の速度を計算すると,次のようになる.
\begin{center}
\begin{tabular}{ccccccccc}
\toprule
$h$ & $1$ & $0.1$ & $0.01$ & $0.001$ & $-0.001$ & $-0.01$ & $-0.1$ & $-1$\\
\midrule
平均の速度 & $5$ & $4.1$ & $4.01$ & $4.001$ & $3.999$ & $3.99$ & $3.9$ & $3$\\
\bottomrule
\end{tabular}
\end{center}
$h$を正の側から$0$に近づけても,負の側から近づけても,平均の速度は$4$に近づいていくように見える.時刻$2$そのものの速度は$4$だと予想できるだろう.

予想を確かめるために,分子を展開してみよう.
\begin{align*}
 \frac{s(2+h)-s(2)}h
 &=\frac{(2+h)^2-2^2}h\\
 &=\frac{4h+h^2}h\\
 &=4+h
\end{align*}
表の値はすべて$4+h$になっていたのである.\,$h$を$0$に近づければ,\,$4+h$は$4$に近づく.

時刻$2$以外ではどうなるだろうか.時刻ごとに同じ計算を繰り返す代わりに,時刻を文字$t$のまま書き下してやろう.
\begin{align*}
 \frac{s(t+h)-s(t)}h
 &=\frac{(t+h)^2-t^2}h\\
 &=\frac{2th+h^2}h\\
 &=2t+h
\end{align*}
$h$を$0$に近づけると$2t$に近づく.したがって,この運動の速度は$v(t)=2t$と表せそうである.文字$t$のまま計算したことで,速度も時刻の関数として得られたのだ.

ただし,\,$h=0$を代入したのではない.\,$h=0$では元の分数の分母が$0$になってしまう.\,$h\ne0$のまま約分し,その後で$h$を$0$に近づけたのである.この「近づける」操作を明確にするのが極限であり,位置から速度を求めた操作を一般の関数にも使えるようにしたものが微分である.
